\section{How common is this design choice?}
\label{sec:generality}

The results so far concern one roster on one board. Two wider questions can be answered without
new labels: whether the choice of verifier also moves the rest of the board, and whether the
scoring design behind the problem is common elsewhere.

\paragraph{Across the whole board.}
The re-scoring, on \BoardRegradeDate{}, rewrote the results file of a model's run with Math-Verify's
scores and did not generate new responses. We traced every displayed score on the board back to its
results file (Appendix~\ref{app:board}). For the \BoardShiftModels{} models whose displayed score
comes from a file the re-scoring rewrote, and whose earlier version we could read, the file's score
before and after the rewrite compares two verifiers on the same responses. Of the pairs that both
scores order strictly, \BoardShiftDiscordantPct\% reverse, that is, swap places
(Appendix~\ref{app:board} defines the statistic and its treatment of ties), and Kendall's rank
correlation between the two scores, which would be $1$ if they ordered every pair alike, is
$\tau_b = \BoardShiftTau$. Among them, on the \RosterShiftModels{} roster models,
\RosterShiftDiscordantPct\% reverse ($\tau_b = \RosterShiftTau$), so the roster understates how far
the choice of verifier moves the board.

\paragraph{Across evaluation frameworks.}
The design choice at the root of the problem is to score an unreadable response as a wrong
answer. We loaded the task configurations of four evaluation frameworks
\citep{biderman2024lessons, opencompass2026, fourrier2024lighteval, openai2024simpleevals} with
each framework's own loader and classified the scoring path of each generative task. In
\texttt{lm-evaluation-harness}, \LmEvalCoerce{} of \LmEvalGenerative{} generative tasks
(\LmEvalCoercePct\%) match an assumed output format and score a failure to match as incorrect. In OpenCompass the figure is
\OpenCompassCoerce{} of \OpenCompassGenerative{} configurations, in \texttt{lighteval}
\LightevalCoerce{} of \LightevalGenerative{} generative metrics, and in \texttt{simple-evals}
\SimpleEvalsCoerce{} of \SimpleEvalsStringMatch{} string-matching modules (Appendix
Table~\ref{tab:frameworks}). A behavioural check, run on a random sample of
\texttt{lm-evaluation-harness} tasks only, confirmed every classification it could test
(Appendix~\ref{app:frameworks}).

These counts say how many tasks could be affected, not how many rankings change; that depends on
how often real responses escape the extraction rule, which we have measured for one benchmark only.
